% Automatisch erzeugt ausanhangsverzeichnis_repo_ano.csv
% Anzahl Einträge: 66

% Benötigte Pakete:
% \usepackage{enumitem}
% \usepackage{xurl}
% \usepackage{microtype} % optional

\begin{enumerate}[label=\arabic*., leftmargin=*, itemsep=0.5em, topsep=0.3em]
\item \textbf{appendix} \textnormal{(\textit{Verzeichnis})}\\
Leerer Ordner, nur zwecks Datenstruktur enthalten.
\item \textbf{code} \textnormal{(\textit{Verzeichnis})}\\
Codes zur Vorverarbeitung der Basisdaten und Berechnung sowie Auswertung der Indikatoren
\begin{enumerate}[label=\arabic{enumi}.\arabic*, leftmargin=*, itemsep=0.5em, topsep=0.3em]
\item \textbf{helper} \textnormal{(\textit{Verzeichnis})}\\
Hilfsfunktionen
\begin{enumerate}[label=\arabic{enumi}.\arabic{enumii}.\arabic*, leftmargin=*, itemsep=0.5em, topsep=0.3em]
\item \textbf{all\_clear.R} \textnormal{(\textit{R})}\\
Stoppen von r5r und Garbage Collection\\
\texttt{\small code/helper/all\_clear.R}
\item \textbf{check\_gtfs\_discont.R} \textnormal{(\textit{R})}\\
Prüfen eines GTFS-Feeds auf Diskontinuitäten\\
\texttt{\small code/helper/check\_gtfs\_discont.R}
\item \textbf{mode\_value.R} \textnormal{(\textit{R})}\\
Berechnen des Modus einer Variable\\
\texttt{\small code/helper/mode\_value.R}
\item \textbf{pointgrid\_fun.R} \textnormal{(\textit{R})}\\
Erstellung einer r5r-kompatiblen Tabelle aus einem Geogitter\\
\texttt{\small code/helper/pointgrid\_fun.R}
\item \textbf{pois\_fun.R} \textnormal{(\textit{R})}\\
Erstellung von r5r-kompatiblen Tabellen für Start- und Zielpunkte\\
\texttt{\small code/helper/pois\_fun.R}
\item \textbf{polygrid\_fun.R} \textnormal{(\textit{R})}\\
Erstellung eines Geogitters\\
\texttt{\small code/helper/polygrid\_fun.R}
\item \textbf{school\_holidays.R} \textnormal{(\textit{R})}\\
Abfrage von Schulferien per API\\
\texttt{\small code/helper/school\_holidays.R}
\end{enumerate}
\item \textbf{temp} \textnormal{(\textit{Verzeichnis})}\\
Verzeichnis zur Ablage temporär erzeugter Dateien (bspw. Untersuchungstage als Output von 03\_find\_valid\_dates
\item \textbf{01\_osm\_extract.R} \textnormal{(\textit{R})}\\
Erzeugen von Grundlagenlayern (Gemeinden, Kreise)
OSM-Zuschnitt auf Untersuchungsgebiet
Vorverarbeitung Zentrales Haltestellenverzeichnis\\
\texttt{\small code/01\_osm\_extract.R}
\item \textbf{02\_de\_gtfs\_cleaning.R} \textnormal{(\textit{R})}\\
Zuschnitt und Vorverarbeitung des DELFI-GTFS\\
\texttt{\small code/02\_de\_gtfs\_cleaning.R}
\item \textbf{03\_find\_valid\_dates.R} \textnormal{(\textit{R})}\\
Identifizieren von validen Analysetagen anhand von Ferien und Feiertagen sowie eventuell fehlenden Daten\\
\texttt{\small code/03\_find\_valid\_dates.R}
\item \textbf{04\_stop\_frequencies\_de\_gtfs.R} \textnormal{(\textit{R})}\\
Berechnung der stündlichen Bedienungsqualitätswerte\\
\texttt{\small code/04\_stop\_frequencies\_de\_gtfs.R}
\item \textbf{05\_bq\_summaries.R} \textnormal{(\textit{R})}\\
Zusammenfassen der BQ nach verschiedenen Zeitfenstern, Erstellung von Plots\\
\texttt{\small code/05\_bq\_summaries.R}
\item \textbf{06\_grid\_creation.R} \textnormal{(\textit{R})}\\
Erstellung eines EPSG:3035-Geogitters mit 100 m Kantenlänge und INSPIRE-konformen Gitter-Ids für das Untersuchungsgebiet\\
\texttt{\small code/06\_grid\_creation.R}
\item \textbf{07\_POI\_creation.R} \textnormal{(\textit{R})}\\
Extraktion und Kategorisierung von POI-Sätzen aus OSM-Daten für Kerndichtenschätzung in 08\_cpt\_KDE.R\\
\texttt{\small code/07\_POI\_creation.R}
\item \textbf{08\_cpt\_KDE.R} \textnormal{(\textit{R})}\\
Kerndichtenschätzung von POI je Gemeinde zur Ermittlung zentralörtlicher Cluster
Bewertung der Cluster und Festlegung zentraler Orte\\
\texttt{\small code/08\_cpt\_KDE.R}
\item \textbf{09\_grid\_zensusvalues.R} \textnormal{(\textit{R})}\\
Verbinden des in 06\_grid\_creation erzeugten Gitters mit den Bevölkerungszahlen\\
\texttt{\small code/09\_grid\_zensusvalues.R}
\item \textbf{10\_i2\_traveltimes.R} \textnormal{(\textit{R})}\\
Berechnung von Gehzeiten (Zensuszellen zu Haltestellen)\\
\texttt{\small code/10\_i2\_traveltimes.R}
\item \textbf{11\_i2\_analysis.R} \textnormal{(\textit{R})}\\
Ermitteln der Erschließungsqualitäten aus Gehzeit und erreichter Bedienungsqualität\\
\texttt{\small code/11\_i2\_analysis.R}
\item \textbf{12\_i2\_summaries.R} \textnormal{(\textit{R})}\\
Aggregieren der EQ-Werte nach Stunde, Tag, Wochentag und Gesamt\\
\texttt{\small code/12\_i2\_summaries.R}
\item \textbf{13\_i2\_togrid.R} \textnormal{(\textit{R})}\\
Erzeugen von Kartenebenen für kartographische Darstellung der EQ, Plot-Darstellungen\\
\texttt{\small code/13\_i2\_togrid.R}
\item \textbf{14\_i3\_traveltimes.R} \textnormal{(\textit{R})}\\
Berechnung von ÖPNV-Wegzeiten (Zentrale Orte zu Zensuszellen)\\
\texttt{\small code/14\_i3\_traveltimes.R}
\item \textbf{15\_i3\_shortest\_times.R} \textnormal{(\textit{R})}\\
Identifizieren kürzester Wegzeiten, Berechnung des Anbindungsindikators je Stunde\\
\texttt{\small code/15\_i3\_shortest\_times.R}
\item \textbf{16\_i3\_summaries.R} \textnormal{(\textit{R})}\\
Aggregieren des Anbindungsindikators nach Tag und Stunde, erzeugen von Kartenebenen und Plots\\
\texttt{\small code/16\_i3\_summaries.R}
\item \textbf{dataenv.R} \textnormal{(\textit{R})}\\
Einträge des Fahrplandatums und ZHV-Datums zur Abfrage in diversen Skripten, Farbpaletten\\
\texttt{\small code/dataenv.R}
\item \textbf{erschließung\_mat\_long\_numeric\_stingy.csv} \textnormal{(\textit{CSV})}\\
Zuordnung von BQ/Gehzeit-Kombinationen zu Erschließungsqualitäten\\
\texttt{\small code/erschließung\_mat\_long\_numeric\_stingy.csv}
\item \textbf{figures\_stop\_frequencies.R} \textnormal{(\textit{R})}\\
Frühere Version von 04\_stop\_frequencies\_de\_gtfs, genutzt um Übersichtsgrafiken zur Bedienungsqualität zu erstellen.\\
\texttt{\small code/figures\_stop\_frequencies.R}
\item \textbf{impedance\_functions.R} \textnormal{(\textit{R})}\\
Erstellung der Grafik zu Impedanzfunktionen in Abschnitt 2.2.1\\
\texttt{\small code/impedance\_functions.R}
\item \textbf{kde\_parallel.R} \textnormal{(\textit{R})}\\
Alternative zu einem Teil von 08\_cpt\_KDE. Parallele Berechnung der KDE-Raster.\\
\texttt{\small code/kde\_parallel.R}
\item \textbf{kreis\_order.txt} \textnormal{(\textit{Text})}\\
Reihenfolge der Kreise nach gewichtetem Mittel der Erschließungsqualität. Output von 13\_i2\_togrid.\\
\texttt{\small code/kreis\_order.txt}
\item \textbf{mean\_frequencies.R} \textnormal{(\textit{R})}\\
Berechnung von Bedienungsqualitäten auf Grundlage täglicher, wochentäglicher und gesamter Mittel der Abfahrten pro Stunde.\\
\texttt{\small code/mean\_frequencies.R}
\item \textbf{osmconf\_cpt.txt} \textnormal{(\textit{Text})}\\
Konfiguration von osmextract zum Extrahieren von POI aus OSM-Daten.\\
\texttt{\small code/osmconf\_cpt.txt}
\item \textbf{stop\_frequencies\_methods.R} \textnormal{(\textit{R})}\\
Frühere Version von 04\_stop\_frequencies\_de\_gtfs, genutzt um Grafiken zur stichprobenartigen Überprüfung der Ergebnisse zu erstellen.
Im weiteren Verlauf nicht aktualisiert, daher nicht funktionsfähig.\\
\texttt{\small code/stop\_frequencies\_methods.R}
\item \textbf{tags\_filter.cmd} \textnormal{(\textit{CMD})}\\
osmium-Ausdruck zum Vorfiltern des OSM-Datensatzes\\
\texttt{\small code/tags\_filter.cmd}
\end{enumerate}
\item \textbf{document} \textnormal{(\textit{Verzeichnis})}\\
\LaTeX-Arbeitsverzeichnis
\begin{enumerate}[label=\arabic{enumi}.\arabic*, leftmargin=*, itemsep=0.5em, topsep=0.3em]
\item \textbf{figures} \textnormal{(\textit{Verzeichnis})}\\
Abbildungen im Dokument
\item \textbf{tables} \textnormal{(\textit{Verzeichnis})}\\
Tabellen im Dokument
\item \textbf{routingindicators.pdf} \textnormal{(\textit{PDF})}\\
Dieses Dokument\\
\texttt{\small document/routingindicators.pdf}
\item \textbf{routingindicators.tex} \textnormal{(\textit{LaTeX})}\\
\TeX-Datei zu diesem Dokument\\
\texttt{\small document/routingindicators.tex}
\end{enumerate}
\item \textbf{feeds} \textnormal{(\textit{Verzeichnis})}\\
Leerer Ordner, nur zwecks Datenstruktur enthalten.
\item \textbf{geodata} \textnormal{(\textit{Verzeichnis})}\\
Geodaten
\begin{enumerate}[label=\arabic{enumi}.\arabic*, leftmargin=*, itemsep=0.5em, topsep=0.3em]
\item \textbf{base\_data} \textnormal{(\textit{Verzeichnis})}\\
Leerer Ordner, nur zwecks Datenstruktur enthalten.
\item \textbf{zhv} \textnormal{(\textit{Verzeichnis})}\\
Leerer Ordner, nur zwecks Datenstruktur enthalten.
\item \textbf{dvg1nw.gpkg} \textnormal{(\textit{GeoPackage})}\\
Verwaltungsgrenzen und Untersuchungsgebiet\\
\texttt{\small geodata/dvg1nw.gpkg}
\begin{enumerate}[label=\arabic{enumi}.\arabic{enumii}.\arabic*, leftmargin=*, itemsep=0.5em, topsep=0.3em]
\item \textbf{regbez10kmbuffer} \textnormal{(\textit{Layer})}\\
Untersuchungsgebiet mit 10 km Puffer (VG250)
\item \textbf{gemeinden\_regbez\_kln} \textnormal{(\textit{Layer})}\\
Gemeindengrenzen des Untersuchungsgebietes (DVG1)
\item \textbf{dvg1krs\_regbez} \textnormal{(\textit{Layer})}\\
Kreisgrenzen des Untersuchungsgebietes (DVG1)
\item \textbf{regbez25kmbuffer} \textnormal{(\textit{Layer})}\\
Untersuchungsgebiet mit 25 km Puffer (VG250)
\item \textbf{gemeinden\_regbez\_25km} \textnormal{(\textit{Layer})}\\
Gemeindengrenzen des Untersuchungsgebietes mit 25 km Puffer (VG250)
\item \textbf{gemeinden\_regbez\_vg250} \textnormal{(\textit{Layer})}\\
Gemeindengrenzen des Untersuchungsgebietes (VG250)
\item \textbf{kreise\_regbez\_vg250} \textnormal{(\textit{Layer})}\\
Kreisgrenzen des Untersuchungsgebietes (VG250)
\end{enumerate}
\item \textbf{regbez10kmbuffer.geojson} \textnormal{(\textit{GeoJSON})}\\
Untersuchungsgebiet mit 10 km Puffer (VG250)\\
\texttt{\small geodata/regbez10kmbuffer.geojson}
\item \textbf{regbez25kmbuffer.geojson} \textnormal{(\textit{GeoJSON})}\\
Untersuchungsgebiet mit 25 km Puffer (VG250)\\
\texttt{\small geodata/regbez25kmbuffer.geojson}
\item \textbf{zensus.gpkg} \textnormal{(\textit{GeoPackage})}\\
Zensusgitter (bewohnte Zellen)\\
\texttt{\small geodata/zensus.gpkg}
\begin{enumerate}[label=\arabic{enumi}.\arabic{enumii}.\arabic*, leftmargin=*, itemsep=0.5em, topsep=0.3em]
\item \textbf{regbez\_zensus\_populated} \textnormal{(\textit{Layer})}
\item \textbf{regbez\_zensus\_populated\_25km} \textnormal{(\textit{Layer})}
\end{enumerate}
\end{enumerate}
\item \textbf{GIS} \textnormal{(\textit{Verzeichnis})}\\
QGIS-Projekt
\begin{enumerate}[label=\arabic{enumi}.\arabic*, leftmargin=*, itemsep=0.5em, topsep=0.3em]
\item \textbf{styles} \textnormal{(\textit{Verzeichnis})}\\
Projektstile
\end{enumerate}
\item \textbf{osmdata} \textnormal{(\textit{Verzeichnis})}\\
Leerer Ordner, nur zwecks Datenstruktur enthalten.
\item \textbf{output} \textnormal{(\textit{Verzeichnis})}\\
Leerer Ordner, nur zwecks Datenstruktur enthalten.
\item \textbf{results} \textnormal{(\textit{Verzeichnis})}\\
Leerer Ordner, nur zwecks Datenstruktur enthalten.
\item \textbf{networkanalysis.Rproj} \textnormal{(\textit{RPROJ})}\\
R-Projektdatei, enthalten für korrekte Pfadinterpretation mit here\\
\texttt{\small networkanalysis.Rproj}
\item \textbf{readme.txt} \textnormal{(\textit{Text})}\\
README\\
\texttt{\small readme.txt}
\end{enumerate}
